% theme: gas_generation_properties
\MajorQuestion{気体の発生と性質}
\SetRowSpaceQanda

\Instruction{次の問いに答えなさい。}
\QQ{空気中に体積の割合で約21％含まれる気体を\NamedBlank{oxygen}という。}{酸素}
\QQ{石灰石にうすい塩酸を加えると発生する気体を\NamedBlank{carbon_dioxide}という。}{二酸化炭素}
\QQ{水に溶けにくい気体を、水と置き換えて集める方法を\NamedBlank{water_displacement}という。}{水上置換法}
\QQ{亜鉛にうすい塩酸を加えると発生する気体を\NamedBlank{hydrogen}という。}{水素}
\QQ{塩化アンモニウムと水酸化カルシウムを混ぜて熱すると発生する気体を\NamedBlank{ammonia}という。}{アンモニア}

\Instruction{\strut}
\QQ{過酸化水素水の市販名を何というか。}{オキシドール}
\QQ{\RefBlank{carbon_dioxide}を通すと、石灰水はどのように変化するか。}{白くにごる}
\QQ{水に溶けやすく、空気より密度が大きい気体を集める方法を\NamedBlank{downward_displacement}という。}{下方置換法}
\QQ{\RefBlank{hydrogen}が燃えるとできる物質は何か。}{水}
\QQ{過酸化水素水に加える黒色の物質を\NamedBlank{manganese_dioxide}という。}{二酸化マンガン}

\Instruction{\strut}
\QQ{空気中に体積の割合で最も多く含まれる気体を\NamedBlank{nitrogen}という。}{窒素}
\QQ{\RefBlank{ammonia}は、水に溶けやすいか溶けにくいか。}{溶けやすい}
\QQ{\RefBlank{carbon_dioxide}が水に少し溶けた水溶液は何性を示すか。}{酸性}
\QQ{水に溶けやすく、空気より密度が小さい気体を集める方法を\NamedBlank{upward_displacement}という。}{上方置換法}
\QQ{\RefBlank{oxygen}がもつ、ものを燃やすはたらきを何というか。}{助燃性}

\Instruction{\strut}
\QQ{燃料の不完全燃焼で生じる、無色・無臭で有毒な気体は何か。}{一酸化炭素}
\QQ{\RefBlank{oxygen}はどの方法で集めるか。}{水上置換法}
\QQ{\RefBlank{carbon_dioxide}は、空気より密度が大きいか小さいか。}{大きい}
\QQ{\RefBlank{ammonia}の水溶液は何性を示すか。}{アルカリ性}
\QQ{\RefBlank{oxygen}は水に溶けやすいか、溶けにくいか。}{溶けにくい}
